\chapter{Conclusion and Outlook}

This thesis set out to answer a deceptively simple question: for a transformer-based classifier
operating on multi-top quark events at the LHC, which architectural choices matter, and why?
Rather than optimising a single model, we built a configurable workbench parameterised along
33 orthogonal axes and used it to run a systematic, controlled ablation programme across
roughly one thousand training runs.

The central finding is that the performance of a transformer classifier in this setting is
determined by a small number of input-representation choices --- principally the tokenisation
strategy, the continuous feature set, and whether physics-informed attention biases are active
--- while the majority of architectural variations imported from the LLM literature, including
normalisation policy, attention-internal normalisation, and causal masking, have no statistically
distinguishable effect on AUROC.
This is itself a practically useful result: it tells a practitioner building a similar
classifier exactly where to invest design effort and where defaults are safe.
The recommended configuration, derived from the axis findings and validated on both the primary
4t-vs-background task and the harder 4t-vs-$t\bar{t}H$ task, achieves a meaningful improvement
in expected discovery significance over the simplest transformer baseline, corresponding to a
reduction in the luminosity required to reach $Z = 5$ that is directly quantifiable in
units of LHC run time.

\section*{Outlook}

Several natural extensions follow from this work.

The most immediate is the incorporation of systematic uncertainties into the sensitivity
estimates of Chapter~10.
The profile likelihood framework used there is already the standard ATLAS machinery; extending
it to include $b$-tagging efficiency, jet energy scale, and luminosity nuisance parameters
would close the gap between the benchmark study presented here and a full physics analysis,
and would test whether the architectural advantages identified in the axes sweep survive the
addition of a realistic uncertainty model.

On the architecture side, the workbench is designed to accept new axes without modifying
existing ones.
Two additions are of particular interest for future work.
First, a graph-based pre-encoder that operates on a dynamically constructed particle graph
--- rather than the fixed-topology MIA blocks studied in Chapter~\ref{ch:preenc} --- would
provide a stronger test of whether explicit geometric structure before the transformer encoder
is beneficial.
Second, the surrogate modelling exercise of Chapter~8 identifies predicted high-performing
configurations in the tails of the explored space; training those configurations would close
the loop between the statistical analysis and the experimental programme.

Beyond the four-top-quark process, the workbench is process-agnostic by construction.
Applying it to three-top production, $t\bar{t}HH$, or other rare multi-object final states at
the HL-LHC requires only a new dataset and a redefinition of the class labels, with all
architectural findings carrying over directly.
The codebase, experiment database, and recommended configuration are publicly available for
exactly this purpose.
